\origpage{164}{ROBERTSON AND SEYMOUR}
Let $\Sigma, \Sigma'$ be connected surfaces, where $\Sigma \simeq \Sigma(a, b, c)$, $\Sigma' \simeq \Sigma(a', b', c')$. We say that $\Sigma$ is \textit{simpler than} $\Sigma'$ if the following conditions are satisfied:

\hspace*{12pt}(i) if $b' = 0$, then $b = 0$;

\hspace*{12pt}(ii) $2a + b \le 2a' + b'$;

\hspace*{12pt}(iii) $4a + 2b + c < 4a' + 2b' + c'$.

From (1.1), this definition does not depend on the choice of $a, b, c, a', b', c'$. \del{Equivalently,} \fix{We use instead the following definition, which implies (i)--(iii) but is not equivalent to them:} $\Sigma$ is simpler than $\Sigma'$ if there exist $a, b, c, a', b', c'$ with $\Sigma \simeq \Sigma(a, b, c)$, $\Sigma' \simeq \Sigma(a', b', c')$, $a \le a'$, $b \le b'$ and $4a + 2b + c < 4a' + 2b' + c'$.\fixnote{fix:1-simpler}{The two definitions are not equivalent. The torus $\Sigma(1, 0, 0)$ and the Klein bottle with one cuff $\Sigma(0, 2, 1)$ satisfy (i)--(iii), but not $a \le a'$. The proof of (13.4), page 207, needs the second form. It uses that a graph that cannot be drawn in $\Sigma'$ cannot be drawn in a surface simpler than $\Sigma'$, and this fails for the first form, since $K_7$ can be drawn in the torus but not in the Klein bottle. With the second form, all the numbered results hold as printed.}

An $O$-arc in a surface $\Sigma$ is a subset homeomorphic to a circle, and a $[0, 1]$-arc is homeomorphic to $[0, 1]$. The ends of a $[0, 1]$-arc are defined in the natural way. An $I$-arc is a $[0, 1]$-arc with both ends in $bd(\Sigma)$ and with no other point in $bd(\Sigma)$.

We shall need to cut along certain $O$- and $I$-arcs. Let $\Sigma$ be a connected surface. If $F$ is an $I$-arc with ends in different cuffs, then we may cut along $F$ and obtain a new surface with one cuff fewer. If $F$ is an $I$-arc with its ends in the same cuff we may cut along $F$ and obtain one of

\hspace*{12pt}(i) a connected surface simpler than $\Sigma$,

\hspace*{12pt}(ii) a surface with two components, both simpler than $\Sigma$, or

\hspace*{12pt}(iii) a surface with two components, one a disc and the other homeomorphic to $\Sigma$.

Finally, if $F \subseteq \Sigma$ is an $O$-arc with $F \cap bd(\Sigma) = \emptyset$, we may cut along $F$ and obtain one of

\hspace*{12pt}(i) a connected surface simpler than $\Sigma$,

\hspace*{12pt}(ii) a surface with two components, both simpler than $\Sigma$,

\hspace*{12pt}(iii) a surface with two components, one a cylinder and the other homeomorphic to $\Sigma$, or

\hspace*{12pt}(iv) a surface with two components, one a disc.

If $X \subseteq \Sigma$, we denote the closure of $X$ by $\bar{X}$, and we denote $\bar{X} - X$ by $\tilde{X}$. A \textit{painting} $\Gamma$ in a surface $\Sigma$ is a pair $(U, N)$, where $U \subseteq \Sigma$ is closed and $N \subseteq U$ is finite, such that

\hspace*{12pt}(i) $U - N$ has only finitely many arc-wise connected components (which we call \textit{cells} of $\Gamma$);

\hspace*{12pt}(ii) for each cell $c$, its closure $\bar{c}$ is a disc and $\tilde{c}$ is a subset of the boundary of this disc, and $|\tilde{c}| \le 3$;

\hspace*{12pt}(iii) $bd(\Sigma) \subseteq U$;

\hspace*{12pt}(iv) for each cell $c$, if $c \cap bd(\Sigma) \ne \emptyset$ then $|\tilde{c}| = 2$ and $\bar{c} \cap bd(\Sigma)$ is a $[0, 1]$-arc with ends the two members of $\tilde{c}$.
